\chapter{Inside and Outside, and What it Means to Assume the Virtue}
\label{chap:pledge-ch4}

%% First use of the epigraph package in the book; set the measure and
%% suppress the default rule so the head sits quietly.
\setlength{\epigraphwidth}{0.62\textwidth}
\setlength{\epigraphrule}{0pt}

\epigraph{\itshape Appearance is the appearance of the interior, and the
interior is the interior of the appearance.}%
{---~Ibn Arabi, \emph{Al-Fut\^uh\^at al-makkiyya}\\(\emph{The Meccan
Illuminations}), vol.~2, p.~427}

\section{Inside and Outside}
\label{sec:pledge-xii}

Once again we find an interesting temporal fold. If you are holding the book in your hand, you have simply turned the page. Meanwhile two months have passed while i have been working out the technical presentation for the Turn. For me, this is where all the juice is. Hence, the lapse between attending to this section of the book and the next. Yet, friends and colleagues who are not in STEM but have read the drafts of those chapters have urged me that i still have much to do by way of bridging between the two cultures, bridging the two sections of the book.

As i have been contemplating how to construct this bridge i have found my attention repeatedly drawn to an understanding that only became vivid for me around about 17 years ago when i was already in my late 40's; and, if i am to be honest, i am still just learning how to navigate this way of apprehending things. Let me see if i can illustrate it in terms of my personal practice.

Everyday, before i allow myself to eat, before i can properly engage the obligations or satisfactions of the day, i do a morning meditation. There is a somewhat involved preparation before i actually settle my ass on the meditation cushion. i won't go into that preparation, except to say that i speak to myself several affirmations and aims and one of them is this: to engage with loving, focused presence. i confess that when the formulation of that affirmation came to me to include in my preparation i was largely thinking of myself. i was asking myself to engage the world and humanity from a place within myself of loving, focused presence.

It is embarrassing to me how many years i stood in front of my cushion, preparing myself to sit before i heard another meaning. It wasn't only about me, it was also a wish to engage Loving Focused Presence as it manifests in the world and the people around me. When i was finally open to this other meaning i realized that the two meanings fit together. Whatever loving focused presence i can muster is a measure of my openness to the Loving Focused Presence, and a measure of what i am willing and able to receive.

The Brits seem to have a handle on this like no other culture. Paul McCartney sang, ``\textbf{[LYRIC TO PASTE: the closing couplet of \emph{The End} --- see the integration notes]}'' And his countryman, Robert Fripp suggests we might ``Assume the virtue.'' i have always taken his words to mean simultaneously: be bold enough to risk being a representative of the virtue; and assume the virtue in others, be generous in your interpretation of their words and deeds. It's a practice that has many dividends, in my experience.

Moreover, it's a message that threads throughout esoteric traditions. By way of example, despite having rubbed shoulders with various Sufis and attended more than one Sema, i only just discovered Ibn Arabi and the wonderful formulation at the head of this chapter.

The reason i mention this wholeness of the interior and the exterior is because i am thinking about my aim in writing this book. i believe that much of the spiritual technology we find articulated and consensually validated across many traditions, such as the one i mentioned above, may be brought into contact with the means of consensus developed in the Western scientific traditions. That crazy Armenian, G.I. Gurdjieff said, ``Take the understanding of the East, and the knowledge of the West --- and then seek.'' What i want to communicate about how we go about recognizing Mind is an attempt to do just that: to bring various ways of understanding and engaging Mind into a discourse where we have the highest quality of rigor. It's in this crucible that genuine, shared understanding is forged.

For example, i'm not interested in any formulation that doesn't admit at least as much precision as a coherent mathematical theory. This is a demand of internal consistency, and it constitutes our culture's highest demand for internal consistency. Of course, lots of sound, coherent mathematical theories have very mysterious connections to physical reality, if any connections at all. If our understanding of Mind is to have real utility --- which in my book is perhaps the best test of all --- it must be physically testable, which ultimately means it must have very concrete connections to physical reality. This is a very good thing! It means that this proposal must also come under the auspices of scientific practice, and scientific rigor.

As i write this i realize that there is another demand. It cannot just be a mathematically rigorous scientific theory. It also has to have a narrative that helps people who do not have a grounding in either maths or science grapple with the ideas, and invites them into the maths and science as crucial means of verifying that our understanding is shared. i believe i said as much at the very outset of this section, in Chapter~\ref{chap:prolegomena}, but it's landing in me differently. At the end of the day, we must each verify for ourselves any given proposal for recognizing Mind. i want us, individually and collectively, to have access to the best tools at our disposal to reach a shared understanding.

To close this off i fear i have painted myself into a corner. Secretly, i began writing this entry to expose a fear that i have become aware of. i am afraid that i will craft a narrative that merely reflects my inner experience. That even if i am able to successfully point to something objective that we can consensually validate together, the framing of the communication will be largely in terms of my inner experience and may, in the end, have less to do with the actual phenomena we are considering than with myself and my personal vision. In some sense i see that this is true. i'm telling a story about what i see with my Mind's eye. i have assiduously selected mathematical presentations that are consistent with what i see. i mentioned the unity of the inside and the outside as a way of reaching toward something that would give me hope that this process also makes an opening for Mind to speak for itself.
